% Fichier complété par tools/complete_missing_corrections.py.
% Source : STAT/STAT-02-Local/50_BancBalafre/50_BancBalafre.tex
% Statut : rédigé (2 question(s)).
\begin{corrigebox}[Corrigé rédigé]
\CorrigeQuestion{1}
Sur l'élément de joint, la pression exerce
                $d\vec F=-p(t)\,dS\,\vec u(\theta)$. Le torseur au point $M$ est
                $\{dT\}_M=\{-p(t)dS\vec u(\theta);\vec0\}_M$.
\CorrigeQuestion{2}
On intègre sur la surface du joint :
                $\vec R=-\int_Sp\vec u\,dS$ et
                $\vec M_B=-\int_S\overrightarrow{BM}\wedge p\vec u\,dS$.
                Les composantes antisymétriques s'annulent; la résultante est portée par
                l'axe de symétrie et vaut la pression multipliée par l'aire projetée.
\end{corrigebox}
